\documentclass[11pt]{article}
\usepackage[a4paper,margin=1in]{geometry}
\usepackage{amsmath,amssymb,amsthm}
\usepackage{graphicx}
\usepackage{xcolor}
\usepackage[hidelinks]{hyperref}
\usepackage{microtype}
\setlength{\emergencystretch}{2em}

\newtheorem{theorem}{Theorem}
\newtheorem{lemma}[theorem]{Lemma}
\newtheorem{corollary}[theorem]{Corollary}
\newcommand{\F}{\mathbb F}
\newcommand{\E}{\mathbb E}
\newcommand{\cH}{\mathcal H}
\newcommand{\supp}{\operatorname{supp}}

\title{Krawtchouk functions maximize every even moment on a Boolean Fourier level}
\author{Automated mathematical discovery run \texttt{fa\_banach\_001}}
\date{13 August 2026}

\begin{document}
\maketitle

\begin{center}
\fcolorbox{black}{gray!8}{\parbox{0.91\linewidth}{
\textbf{Status.} Candidate full solution, likely valid, subject to human
review.  For every even exponent $p=2m$, the two maxima in the source's
appendix Question B.2(6) are attained by a permutation-invariant
Krawtchouk function.  The $p=4$ case was proved by Aaronson in 2018; the
argument below extends the compression principle to all even moments.
}}
\end{center}

\section{Source question and provenance}

The source is Y.~Polyanskiy, \emph{Hypercontractivity of spherical averages
in Hamming space}, arXiv:1309.3014 \cite{Source}.  Write
$G=\F_2^n$, equip it with uniform probability measure, and let
$\Pi_a$ denote orthogonal projection onto the Fourier level

\[
 \cH_a=\operatorname{span}\{\chi_A:A\subseteq[n],\ |A|=a\},
 \qquad \chi_A(x)=(-1)^{\sum_{i\in A}x_i}.
\]

The source asks whether, for $p=2m$, the maxima of

\[
 \frac{\|f_a\|_p}{\|f_a\|_2}
 \quad\text{and}\quad
 \frac{\|f_a\|_p}{\|f\|_2},\qquad f_a=\Pi_a f,
\]

are achieved by permutation-invariant functions.  The exact passage is
reproduced below.

\begin{center}
\includegraphics[width=0.96\linewidth]{figures/open_problem_crop.png}
\end{center}

\noindent\textbf{Provenance note.}
This question occurs in the arXiv source tree inside an appendix guarded by
the source's \texttt{ifmapx} switch; it is not printed in the
official 24-page arXiv PDF.  The crop is from page 27 of a DVI/PDF build of
the original source with its documented appendix job name
\texttt{hc\_hamming\_apx}.  The official PDF is retained as
\texttt{source\_paper.pdf}.

\section{Result and proof intuition}

Define the degree-$a$ Krawtchouk function
\[
 K_a(x)=\sum_{|A|=a}\chi_A(x),
 \qquad \|K_a\|_2=\binom na^{1/2}.
\]

\begin{theorem}[Exact even-moment level norm]\label{thm:main}
For all $n\geq1$, $0\leq a\leq n$, and integers $m\geq1$,
\begin{equation}\label{eq:level}
 \sup_{0\ne g\in\cH_a}\frac{\|g\|_{2m}}{\|g\|_2}
 =\frac{\|K_a\|_{2m}}{\|K_a\|_2}.
\end{equation}
The supremum may be taken over real- or complex-valued functions.  Moreover,
\begin{equation}\label{eq:projection}
 \|\Pi_a\|_{L_2(G)\to L_{2m}(G)}
 =\frac{\|K_a\|_{2m}}{\|K_a\|_2}.
\end{equation}
Thus both maxima in the source question are attained by the
permutation-invariant function $K_a$.
\end{theorem}

The proof works on the Fourier-coefficient sphere
$S(n,a)=\{A\subseteq[n]:|A|=a\}$.  First discard coefficient phases: the
even-moment expansion can only increase when every coefficient is replaced
by its modulus.  Next choose two coordinates $i,j$.  Whenever two
coefficients sit at the endpoints of a Johnson-graph edge
$A\leftrightarrow A\mathbin\triangle\{i,j\}$, replace both by their
root-mean-square.  This preserves the coefficient $\ell_2$ norm.

For the fourth moment, monotonicity of this compression is Aaronson's key
observation \cite{Aaronson}.  At a general $2m$-th moment, fix all
coordinates outside $i,j$ in the $2m$-tuple expansion.  The remaining
sum is an even- or odd-parity tensor.  Its elementary sharp product bound
shows that the same root-mean-square compression is monotone for every
even moment.  Repeating the compressions averages the squared coefficients
over the connected Johnson graph, leaving the constant vector and hence
$K_a$.

\section{The parity-tensor inequality}

\begin{lemma}[Parity block bound]\label{lem:parity}
Let $r\geq2$, let $u_\ell,v_\ell\geq0$, and let
$\epsilon\in\{0,1\}$.  Then
\begin{equation}\label{eq:parity}
 \sum_{\substack{\sigma\in\{0,1\}^r\\
                  \sum_\ell\sigma_\ell\equiv\epsilon\ (2)}}
 \prod_{\ell=1}^r
 \bigl(u_\ell\mathbf1_{\{\sigma_\ell=0\}}
      +v_\ell\mathbf1_{\{\sigma_\ell=1\}}\bigr)
 \leq
 2^{r/2-1}\prod_{\ell=1}^r(u_\ell^2+v_\ell^2)^{1/2}.
\end{equation}
The right side is exactly the left side after every pair
$(u_\ell,v_\ell)$ is replaced by
$(\rho_\ell,\rho_\ell)$, where
$\rho_\ell=((u_\ell^2+v_\ell^2)/2)^{1/2}$.
\end{lemma}

\begin{proof}
The parity filter writes the left side as
\[
 \frac12\left\{\prod_\ell(u_\ell+v_\ell)
 +(-1)^\epsilon\prod_\ell(u_\ell-v_\ell)\right\}.
\]
Zero pairs may be removed, so normalize
$t_\ell=(u_\ell^2+v_\ell^2)^{1/2}>0$, and set
\[
 x_\ell=\frac{u_\ell+v_\ell}{\sqrt2,t_\ell},\qquad
 y_\ell=\frac{u_\ell-v_\ell}{\sqrt2,t_\ell}.
\]
Then $x_\ell^2+y_\ell^2=1$.  Since $r\geq2$,
\[
 \prod_\ell|x_\ell|+\prod_\ell|y_\ell|
 \leq |x_1x_2|+|y_1y_2|
 \leq\sqrt{(x_1^2+y_1^2)(x_2^2+y_2^2)}=1.
\]
Substitution proves \eqref{eq:parity}.  After replacement there are
$2^{r-1}$ admissible parity strings, each contributing
$2^{-r/2}\prod t_\ell$, proving the final assertion.
\end{proof}

\section{Proof of the theorem}

\begin{proof}[Proof of Theorem~\ref{thm:main}]
Write $g=\sum_{|A|=a}c_A\chi_A$.  Expanding against the orthonormal
characters gives
\begin{equation}\label{eq:moment}
 \|g\|_{2m}^{2m}
 =\sum_{\substack{A_1,\ldots,A_m,B_1,\ldots,B_m\in S(n,a)\\
 A_1\triangle\cdots\triangle A_m
 =B_1\triangle\cdots\triangle B_m}}
 c_{A_1}\cdots c_{A_m}
 \overline{c_{B_1}\cdots c_{B_m}}.
\end{equation}
Taking absolute values term by term shows that this is at most the same
moment with $c_A$ replaced by $|c_A|$, while Parseval leaves
$\|g\|_2^2=\sum_A|c_A|^2$ unchanged.  It therefore suffices to assume
all coefficients are nonnegative.

Fix $i<j$.  If $A$ contains exactly one of $i,j$, put
$A'=A\triangle\{i,j\}$ and replace both $c_A,c_{A'}$ by
\[
 \rho_A=\left(\frac{c_A^2+c_{A'}^2}{2}\right)^{1/2}.
\]
Leave the other coefficients fixed.  This compression preserves
$\sum_Ac_A^2$.

To compare moments, use the equivalent $2m$-tuple form of
\eqref{eq:moment}, whose constraint is
$A_1\triangle\cdots\triangle A_{2m}=\varnothing$.  Group terms after
deleting coordinates $i,j$ from every $A_\ell$.  If $r$ of the
resulting bases have size $a-1$, their lifts have coordinate patterns
$10$ or $01$.  The other lifts, of sizes $a$ and $a-2$, are fixed
as $00$ and $11$.  A block contributes only when $r$ is even.  For
$r=0$ it is unchanged.  For even $r\geq2$, the two coordinate parity
conditions reduce to one prescribed parity of the $r$ binary choices.
After removing the common product of fixed coefficients, the block is the
left side of \eqref{eq:parity}.  Lemma~\ref{lem:parity} says precisely that
the compression cannot decrease it.  Hence every coordinate-pair
compression preserves $\|g\|_2$ and increases $\|g\|_{2m}$.

On squared coefficients $w_A=c_A^2$, this operation is ordinary averaging
across every Johnson-graph edge carrying the label $i,j$.  Cycle through
all coordinate pairs and repeat.  These doubly stochastic averaging maps
have a positive diagonal, and the union of their edges is the connected
Johnson graph $J(n,a)$.  The standard finite consensus argument therefore
gives
\[
 w_A\longrightarrow \binom na^{-1}\sum_{|B|=a}w_B
 \qquad\text{for every }A\in S(n,a).
\]
Continuity of the finite moment polynomial shows that the original ratio is
at most the ratio for constant nonnegative coefficients.  Their Fourier
extension is a scalar multiple of $K_a$, proving \eqref{eq:level}.

Finally, for arbitrary $f$, orthogonality gives
\[
 \|\Pi_af\|_{2m}
 \leq \frac{\|K_a\|_{2m}}{\|K_a\|_2}\|\Pi_af\|_2
 \leq \frac{\|K_a\|_{2m}}{\|K_a\|_2}\|f\|_2.
\]
Equality is attained at $f=K_a$, which proves \eqref{eq:projection} and
the second assertion in the source question.
\end{proof}

\section{Verification, literature boundary, and limitations}

The proof was checked in three independent ways.  First, direct numerical
optimization over every tested level through $n=7$, for moments
$p=4,6,8$, returned the constant-magnitude coefficient orbit.  Second,
20,000 random pair-compression tests in each of nine configurations through
$p=10$ found no decrease.  Third, the block formula was checked directly
against enumeration of parity strings.  These computations are sanity
checks only; the proof is the exact parity-tensor argument above.

The case $p=4$ is not new: Aaronson proved exact constant-coefficient
extremality for additive energy \cite{Aaronson}, answering a formulation of
Kirshner--Samorodnitsky \cite{KS18}.  Kirshner--Samorodnitsky later obtained
general-$p$ moment bounds that are asymptotically tight up to lower-order
factors \cite{KS19}; that work does not state the exact all-even-moment
identity above.  Bounded searches through 13 August 2026 used the exact
source wording, ``Krawtchouk extremizer'', ``even moments'', higher additive
energy, and the citation neighborhood of arXiv:1805.05295.  No exact
all-even-$p$ theorem was found.  Novelty confidence is therefore moderate,
not definitive.

The theorem addresses only the even exponents explicitly proposed in the
source.  It makes no claim for non-even $p>2$, where the polynomial
expansion and parity compression are unavailable; the source already notes
failure of the analogous symmetry assertion at degree one for
$2<p<3$.  It also does not classify all equality cases.

\begin{thebibliography}{9}
\bibitem{Source}
Y.~Polyanskiy,
\emph{Hypercontractivity of spherical averages in Hamming space},
arXiv:1309.3014.

\bibitem{Aaronson}
J.~Aaronson,
\emph{Functions with large additive energy supported on a Hamming sphere},
arXiv:1805.05295, 2018.

\bibitem{KS18}
N.~Kirshner and A.~Samorodnitsky,
\emph{On the $\ell_4:\ell_2$ ratio of functions with restricted Fourier
support}, arXiv:1801.08507.

\bibitem{KS19}
N.~Kirshner and A.~Samorodnitsky,
\emph{A moment ratio bound for polynomials and some extremal properties of
Krawtchouk polynomials and Hamming spheres}, arXiv:1909.11929.
\end{thebibliography}

\end{document}
